\begin{theorem}
  \label{ent:thm:linear-modeq-iff}
  For integers $n,a,b,x$, the congruence $ax\equiv b\pmod n$ is equivalent to the existence of an
  integer $y$ with $ax+ny=b$.

  \begin{lean}
    theorem linear_modEq_iff (n a b x : ℤ) :
        a * x ≡ b [ZMOD n] ↔ ∃ y : ℤ, a * x + n * y = b
  \end{lean}
\end{theorem}

\begin{proof}
  Negate the witness for $n\mid ax-b$ to obtain $y$, and reverse this construction in the other
  direction.

  \begin{lean}
    rw [modEq_iff_dvd_sub]
    constructor

    · rintro ⟨k, hk⟩
      exact ⟨-k, by nlinarith⟩

    · rintro ⟨y, hy⟩
      exact ⟨-y, by nlinarith⟩
  \end{lean}
\end{proof}

\begin{theorem}
  \label{ent:thm:exists-linear-modeq-iff}
  The congruence $ax\equiv b\pmod n$ is solvable exactly when $\gcd(a,n)\mid b$.

  \begin{lean}
    theorem exists_linear_modEq_iff (n a b : ℤ) :
        (∃ x : ℤ, a * x ≡ b [ZMOD n]) ↔ (Chapter01.gcd a n : ℤ) ∣ b
  \end{lean}
\end{theorem}

\begin{proof}
  Use the preceding translation and the solvability criterion for linear Diophantine equations.

  \begin{lean}
    simp only [linear_modEq_iff]
    exact Chapter01.exists_mul_add_mul_eq_iff a n b
  \end{lean}
\end{proof}

\begin{theorem}
  \label{ent:thm:linear-modeq-iff-eq-add}
  Let $n\ne0$, $g=\gcd(a,n)$, and let $x_0$ solve $ax\equiv b\pmod n$.
  The solutions are precisely $x=x_0+(n/g)t$ with $t\in\mathbb Z$.

  \begin{lean}
    theorem linear_modEq_iff_eq_add (n a b x₀ x : ℤ) (hn : n ≠ 0)
        (h₀ : a * x₀ ≡ b [ZMOD n]) :
        a * x ≡ b [ZMOD n] ↔ ∃ t : ℤ,
          x = x₀ + (n / (Chapter01.gcd a n : ℤ)) * t
  \end{lean}
\end{theorem}

\begin{proof}
  Choose a companion $y_0$ with $ax_0+ny_0=b$.
  For any other solution, the previously proved parametrization of Diophantine solutions gives the
  displayed first coordinate.

  \begin{lean}
    obtain ⟨y₀, hy₀⟩ := (linear_modEq_iff n a b x₀).1 h₀
    rw [linear_modEq_iff]
    constructor

    · rintro ⟨y, hy⟩

      obtain ⟨t, ht, _⟩ :=
        (Chapter01.mul_add_mul_eq_iff a n b x₀ y₀ x y (Or.inr hn) hy₀).1 hy

      exact ⟨t, ht⟩
  \end{lean}

  Conversely, accompany any such $x$ by $y_0-(a/g)t$, where $g=\gcd(a,n)$.
  The parametrization then verifies the original equation.

  \begin{lean}
    · rintro ⟨t, ht⟩
      refine ⟨y₀ - (a / (Chapter01.gcd a n : ℤ)) * t, ?_⟩
      exact (Chapter01.mul_add_mul_eq_iff a n b x₀ y₀ x _ (Or.inr hn) hy₀).2
        ⟨t, ht, rfl⟩
  \end{lean}
\end{proof}

\begin{theorem}
  \label{ent:thm:linear-solution-classes}
  For $n>0$, a solvable linear congruence has exactly $g=\gcd(a,n)$ incongruent solutions.
  They are represented by $x_0+(n/g)t$ for $0\le t<g$, and each solution belongs to exactly one of
  these classes.

  \begin{lean}
    theorem linear_solution_classes (n a b x₀ : ℤ) (hn : 0 < n)
        (h₀ : a * x₀ ≡ b [ZMOD n]) :
        (∀ t : Fin (Chapter01.gcd a n),
          a * (x₀ + (n / (Chapter01.gcd a n : ℤ)) * t.val) ≡ b [ZMOD n]) ∧
        (∀ x : ℤ, a * x ≡ b [ZMOD n] →
          ∃! t : Fin (Chapter01.gcd a n),
            x ≡ x₀ + (n / (Chapter01.gcd a n : ℤ)) * t.val [ZMOD n])
  \end{lean}
\end{theorem}

\begin{proof}
  Write $n=gu$ with $u>0$.
  Every proposed representative is a solution by the parametrization.

  \begin{lean}
    have hnne := ne_of_gt hn

    obtain ⟨hg, _, ⟨u, hu⟩, _, _⟩ := Chapter01.gcd_spec a n (Or.inr hnne)

    have hgne : (Chapter01.gcd a n : ℤ) ≠ 0 := by
      exact_mod_cast ne_of_gt hg

    have hquot : n / (Chapter01.gcd a n : ℤ) = u := by
      calc
        n / (Chapter01.gcd a n : ℤ) =
            ((Chapter01.gcd a n : ℤ) * u) / (Chapter01.gcd a n : ℤ) :=
          congrArg (fun z => z / (Chapter01.gcd a n : ℤ)) hu
        _ = u := Int.mul_ediv_cancel_left u hgne

    have hupos : 0 < u := by
      have hgz : (0 : ℤ) < Chapter01.gcd a n := by
        exact_mod_cast hg
      nlinarith

    constructor

    · intro t
      exact (linear_modEq_iff_eq_add n a b x₀ _ hnne h₀).2 ⟨t.val, rfl⟩
  \end{lean}

  For an arbitrary solution, divide its integer parameter by $g=\gcd(a,n)$ and retain the
  remainder $r$.
  The discarded multiple of $g$ changes $x$ by a multiple of $n$, so $r$ supplies a representative.

  \begin{lean}
    · intro x hx

      obtain ⟨t, ht⟩ := (linear_modEq_iff_eq_add n a b x₀ x hnne h₀).1 hx
      obtain ⟨⟨q, r⟩, ⟨htr, hr⟩, _⟩ :=
        Chapter01.existsUnique_quotient_remainder t ⟨Chapter01.gcd a n, hg⟩

      change t = q * (Chapter01.gcd a n : ℤ) + r at htr
      change r < Chapter01.gcd a n at hr

      have hrep : x ≡ x₀ + (n / (Chapter01.gcd a n : ℤ)) * r [ZMOD n] := by
        apply (modEq_iff_dvd_sub _ _ _).2
        refine ⟨q, ?_⟩
        rw [hquot] at ht ⊢
        rw [ht]
        nlinarith [congrArg (fun z => z * q) hu, congrArg (fun z => u * z) htr]
  \end{lean}

  If another remainder $s$ represents the same class, subtracting the two congruences and
  cancelling $u>0$ makes $s-r$ a multiple of $g$.

  \begin{lean}
    refine ⟨⟨r, hr⟩, hrep, ?_⟩
    intro s hs

    have hdiff := modEq_trans (modEq_symm hs) hrep

    obtain ⟨k, hk⟩ := (modEq_iff_dvd_sub _ _ _).1 hdiff
    rw [hquot] at hk

    have he : (s.val : ℤ) - r = (Chapter01.gcd a n : ℤ) * k := by
      apply mul_left_cancel₀ (ne_of_gt hupos)
      calc
        u * ((s.val : ℤ) - r) = (x₀ + u * s.val) - (x₀ + u * r) := by ring
        _ = u * ((Chapter01.gcd a n : ℤ) * k) := by
          rw [hk]
          nlinarith [congrArg (fun z => z * k) hu]
  \end{lean}

  Both remainders lie between zero and $g-1$.
  Their difference is therefore strictly between $-g$ and $g$, forcing that multiple to be zero.

  \begin{lean}
    have hsbound : (s.val : ℤ) < Chapter01.gcd a n := by
      exact_mod_cast s.isLt

    have hrbound : (r : ℤ) < Chapter01.gcd a n := by
      exact_mod_cast hr

    have hgz : (0 : ℤ) < Chapter01.gcd a n := by
      exact_mod_cast hg

    have hkzero : k = 0 := by
      rcases lt_trichotomy k 0 with hneg | hzero | hpos

      · have : k ≤ -1 := by
          omega
        nlinarith [Int.natCast_nonneg s.val, Int.natCast_nonneg r]

      · exact hzero

      · have : 1 ≤ k := hpos
        nlinarith [Int.natCast_nonneg s.val, Int.natCast_nonneg r]
  \end{lean}

  Thus $s=r$, proving uniqueness of the representative.

  \begin{lean}
    apply Fin.ext

    have hval : (s.val : ℤ) = r := by
      rw [hkzero, mul_zero] at he
      linarith
    exact_mod_cast hval
  \end{lean}
\end{proof}

\begin{definition}
  An integer $u$ is a multiplicative inverse of $a$ modulo $n$ when $au\equiv1\pmod n$.

  \begin{lean}
    def IsModularInverse (n a u : ℤ) : Prop := a * u ≡ 1 [ZMOD n]
  \end{lean}
\end{definition}

\begin{theorem}
  \label{ent:thm:exists-modularinverse}
  If $\gcd(a,n)=1$, an inverse of $a$ modulo $n$ exists.

  \begin{lean}
    theorem exists_modularInverse (n a : ℤ) (ha : Chapter01.gcd a n = 1) :
        ∃ u : ℤ, IsModularInverse n a u
  \end{lean}
\end{theorem}

\begin{proof}
  Apply the linear congruence existence criterion with right-hand side one.

  \begin{lean}
    apply (exists_linear_modEq_iff n a 1).2
    rw [ha, Nat.cast_one]
  \end{lean}
\end{proof}

\begin{theorem}
  \label{ent:thm:exists-unique-linear-modeq}
  For $n\ne0$ and $\gcd(a,n)=1$, every linear congruence $ax\equiv b\pmod n$ has one solution
  class modulo $n$.

  \begin{lean}
    theorem exists_unique_linear_modEq (n a b : ℤ) (hn : n ≠ 0)
        (ha : Chapter01.gcd a n = 1) :
        ∃ x₀ : ℤ, a * x₀ ≡ b [ZMOD n] ∧
          ∀ x : ℤ, a * x ≡ b [ZMOD n] ↔ x ≡ x₀ [ZMOD n]
  \end{lean}
\end{theorem}

\begin{proof}
  Existence follows because one divides every integer.

  \begin{lean}
    obtain ⟨x₀, hx₀⟩ := (exists_linear_modEq_iff n a b).2 (by
      rw [ha, Nat.cast_one]
      exact ⟨b, by ring⟩)
  \end{lean}

  The solution step $n/\gcd(a,n)$ is now $n$, so the parametrization is exactly one congruence
  class.

  \begin{lean}
    refine ⟨x₀, hx₀, ?_⟩
    intro x
    rw [linear_modEq_iff_eq_add n a b x₀ x hn hx₀, ha, Nat.cast_one, Int.ediv_one,
      modEq_iff_dvd_sub]

    constructor

    · rintro ⟨t, ht⟩
      exact ⟨t, by linarith⟩

    · rintro ⟨t, ht⟩
      exact ⟨t, by linarith⟩
  \end{lean}
\end{proof}

\begin{theorem}
  \label{ent:thm:linear-modeq-iff-inverse}
  If $au\equiv1\pmod n$, then $ax\equiv b\pmod n$ is equivalent to $x\equiv ub\pmod n$.

  \begin{lean}
    theorem linear_modEq_iff_inverse (n a b u x : ℤ) (hu : IsModularInverse n a u) :
        a * x ≡ b [ZMOD n] ↔ x ≡ u * b [ZMOD n]
  \end{lean}
\end{theorem}

\begin{proof}
  Suppose $ax\equiv b\pmod n$.
  Multiplication by $u$ gives $u(ax)\equiv ub$, while the inverse relation gives $(au)x\equiv x$.

  \begin{lean}
    constructor

    · intro hx

      have h₁ := modEq_mul (modEq_refl n u) hx
      have h₂ := modEq_mul hu (modEq_refl n x)
  \end{lean}

  Since $u(ax)=(au)x$, these congruences imply $x\equiv ub\pmod n$.

  \begin{lean}
    have heq : u * (a * x) = (a * u) * x := by
      ring
    rw [heq] at h₁

    exact modEq_trans (modEq_symm (by simpa only [one_mul] using h₂)) h₁
  \end{lean}

  Conversely, suppose $x\equiv ub\pmod n$.
  Multiplication by $a$ gives $ax\equiv a(ub)$, and the inverse relation gives $(au)b\equiv b$.

  \begin{lean}
    · intro hx

      have h₁ := modEq_mul (modEq_refl n a) hx
      have h₂ := modEq_mul hu (modEq_refl n b)
  \end{lean}

  Rewriting $a(ub)=(au)b$ yields $ax\equiv b\pmod n$.

  \begin{lean}
    have heq : a * (u * b) = (a * u) * b := by
      ring
    rw [heq] at h₁

    exact modEq_trans h₁ (by simpa only [one_mul] using h₂)
  \end{lean}
\end{proof}

\begin{theorem}
  \label{ent:thm:chinese-remainder}
  Let $(n_i)_{i\in S}$ be a finite family of nonzero, pairwise coprime integer moduli.
  For any integers $a_i$, the system $x\equiv a_i\pmod{n_i}$ has a solution, and all solutions form
  one class modulo $\prod_{i\in S}n_i$.
  The empty system gives the unique class modulo one.

  \begin{lean}
    theorem chinese_remainder {ι : Type*} (s : Finset ι) (n a : ι → ℤ)
        (hn : ∀ i ∈ s, n i ≠ 0)
        (hpair : ∀ i ∈ s, ∀ j ∈ s, i ≠ j → Chapter01.gcd (n i) (n j) = 1) :
        ∃ x : ℤ, (∀ i ∈ s, x ≡ a i [ZMOD n i]) ∧
          ∀ y : ℤ, (∀ i ∈ s, y ≡ a i [ZMOD n i]) ↔ y ≡ x [ZMOD ∏ i ∈ s, n i]
  \end{lean}
\end{theorem}

\begin{proof}
  For each $i$, form the complementary product $N_i=\prod_{j\ne i}n_j$.
  Pairwise coprimality and the product form of Bezout\textquotesingle s identity give an inverse
  $u_i$ of $N_i$ modulo $n_i$.
  Set $x=\sum_i a_iN_i u_i$.

  \begin{lean}
    classical

    let N (i : ι) := ∏ j ∈ s.erase i, n j

    have hcop (i : ι) (hi : i ∈ s) : Chapter01.gcd (N i) (n i) = 1 := by
      apply gcd_prod_eq_one _ _ _ (hn i hi)
      intro j hj
      exact hpair j (Finset.mem_of_mem_erase hj) i hi (Finset.ne_of_mem_erase hj)

    have hinv (i : ι) (hi : i ∈ s) : ∃ u : ℤ, IsModularInverse (n i) (N i) u :=
      exists_modularInverse _ _ (hcop i hi)
    choose u hu using hinv

    let x := ∑ i ∈ s.attach, a i * N i * u i i.property
  \end{lean}

  Fix $i$.
  Subtracting $a_i$ from the constructed sum amounts to subtracting $a_i$ from its $i$th summand
  and leaving the others unchanged.

  \begin{lean}
    have hx (i : ι) (hi : i ∈ s) : x ≡ a i [ZMOD n i] := by
      have hsum : x - a i =
          ∑ j ∈ s.attach, (a j * N j * u j j.property - if j.val = i then a i else 0) := by
        simp only [Finset.sum_sub_distrib]

        have he : (∑ j ∈ s.attach, if j.val = i then a i else 0) = a i := by
          rw [Finset.sum_eq_single ⟨i, hi⟩]

          · simp only [↓reduceIte]

          · intro j hj hji

            have hne : j.val ≠ i := by
              intro heq
              apply hji
              exact Subtype.ext heq
            simp only [hne, ↓reduceIte]

          · simp
        rw [he]
  \end{lean}

  For the $i$th summand, the inverse relation makes the adjusted term divisible by $n_i$.

  \begin{lean}
    apply (modEq_iff_dvd_sub _ _ _).2
    rw [hsum]
    apply Finset.dvd_sum
    intro j hj
    by_cases he : j.val = i

    · subst i
      simp only [↓reduceIte]

      obtain ⟨k, hk⟩ := (modEq_iff_dvd_sub _ _ _).1 (hu j.val j.property)

      exact ⟨a j * k, by nlinarith [congrArg (fun z => a j * z) hk]⟩
  \end{lean}

  Every other summand contains $n_i$ in its complementary product.
  Each adjusted summand is thus divisible by $n_i$, so their sum is too.

  \begin{lean}
    · simp only [he, ↓reduceIte, sub_zero]

      have hmem : i ∈ s.erase j.val := Finset.mem_erase.mpr ⟨Ne.symm he, hi⟩

      obtain ⟨k, hk⟩ := Finset.dvd_prod_of_mem n hmem

      exact ⟨a j * k * u j j.property, by
          change N j = n i * k at hk
          rw [hk]
          ring⟩
  \end{lean}

  Two simultaneous solutions have a difference divisible by every $n_i$.
  Repeated use of the coprime product-divisibility lemma makes the difference divisible by their
  product.
  Conversely, a multiple of that product is a multiple of each modulus.

  \begin{lean}
    refine ⟨x, hx, ?_⟩
    intro y
    constructor

    · intro hy
      apply (modEq_iff_dvd_sub _ _ _).2
      apply prod_dvd_of_pairwise s n hn hpair
      intro i hi
      exact (modEq_iff_dvd_sub _ _ _).1 (modEq_trans (hy i hi) (modEq_symm (hx i hi)))

    · intro hy i hi

      obtain ⟨k, hk⟩ := (modEq_iff_dvd_sub _ _ _).1 hy
      obtain ⟨t, ht⟩ := Finset.dvd_prod_of_mem n hi
      apply modEq_trans _ (hx i hi)
      exact (modEq_iff_dvd_sub _ _ _).2 ⟨t * k, by
          rw [hk, ht]
          ring⟩
  \end{lean}
\end{proof}

\begin{theorem}
  \label{ent:thm:linear-system-iff}
  Let $\Delta=ad-bc$ and suppose $\Delta t\equiv1\pmod n$.
  The system $ax+by\equiv r$, $cx+dy\equiv s\pmod n$ is equivalent to $x\equiv t(dr-bs)$ and
  $y\equiv t(as-cr)\pmod n$.

  \begin{lean}
    theorem linear_system_iff (n : ℕ) (a b c d r s t x y : ℤ)
        (ht : (a * d - b * c) * t ≡ 1 [ZMOD (n : ℤ)]) :
        (a * x + b * y ≡ r [ZMOD (n : ℤ)] ∧ c * x + d * y ≡ s [ZMOD (n : ℤ)]) ↔
          x ≡ t * (d * r - b * s) [ZMOD (n : ℤ)] ∧
          y ≡ t * (a * s - c * r) [ZMOD (n : ℤ)]
  \end{lean}
\end{theorem}

\begin{proof}
  Work in the residue ring.
  Subtract $b$ times the second equation from $d$ times the first, and subtract $c$ times the first
  from $a$ times the second.
  These give $\Delta x=dr-bs$ and $\Delta y=as-cr$.
  Multiplication by $t$ gives the formulas; direct substitution verifies their converse.
  The auxiliary ring calculation is proved in the same checked module.

  \begin{lean}
    simp only [← ZMod.intCast_eq_intCast_iff, Int.cast_add, Int.cast_mul,
      Int.cast_sub, Int.cast_one] at ht ⊢

    exact two_equations (a : ZMod n) b c d r s t x y ht
  \end{lean}
\end{proof}

\begin{theorem}
  \label{ent:thm:exists-unique-linear-system}
  If $\gcd(ad-bc,n)=1$, the two-variable system has a unique solution pair modulo $n$.

  \begin{lean}
    theorem exists_unique_linear_system (n : ℕ) (a b c d r s : ℤ)
        (hdet : Chapter01.gcd (a * d - b * c) n = 1) :
        ∃ x₀ y₀ : ℤ, ∀ x y : ℤ,
          (a * x + b * y ≡ r [ZMOD (n : ℤ)] ∧ c * x + d * y ≡ s [ZMOD (n : ℤ)]) ↔
            x ≡ x₀ [ZMOD (n : ℤ)] ∧ y ≡ y₀ [ZMOD (n : ℤ)]
  \end{lean}
\end{theorem}

\begin{proof}
  Choose an inverse of the determinant and apply the explicit solution formulas.

  \begin{lean}
    obtain ⟨t, ht⟩ := exists_modularInverse n (a * d - b * c) hdet

    exact ⟨t * (d * r - b * s), t * (a * s - c * r),
      fun x y => linear_system_iff n a b c d r s t x y ht⟩
  \end{lean}
\end{proof}
